\documentclass[UTF8,fontset=windows]{ctexart}

\usepackage[a4paper,margin=2.4cm]{geometry}
\usepackage{booktabs}
\usepackage{enumitem}
\usepackage{xcolor}
\usepackage{hyperref}
\usepackage{fvextra}

\hypersetup{
  colorlinks=true,
  linkcolor=blue,
  urlcolor=blue
}

\setlist{nosep}
\DefineVerbatimEnvironment{code}{Verbatim}{
  fontsize=\footnotesize,
  breaklines=true,
  breakanywhere=true,
  frame=single,
  rulecolor=\color{black!25}
}

\title{FG-MEET 程序执行方案}
\author{}
\date{2026 年 6 月 25 日}

\begin{document}

\maketitle

\section{执行目标}

本文档用于说明 FG-MEET 交付包在 MATLAB 中的正确执行方式。执行流程包括：

\begin{enumerate}
  \item 进入交付包根目录；
  \item 初始化 MATLAB 路径；
  \item 运行快速验证算例；
  \item 按需指定单个 \texttt{txt} 输入文件执行力、电、磁三类工况；
  \item 必要时重新生成功能梯度 \texttt{txt} 输入文件。
\end{enumerate}

主要求解入口是项目根目录下的脚本和 \texttt{run\_meet\_static} 函数，不是 \texttt{materials} 文件夹中的材料辅助函数。

\section{目录要求}

首先进入交付包根目录。该目录应包含以下文件和文件夹：

\begin{itemize}
  \item \texttt{setup\_paths.m}
  \item \texttt{run\_phase1\_static\_elastic.m}
  \item \texttt{run\_meet\_static.m}
  \item \texttt{matlab/}
  \item \texttt{materials/}
  \item \texttt{cases/}
  \item \texttt{tools/}
\end{itemize}

示例路径如下：

\begin{code}
C:\Users\You Huiwen\Desktop\fg-meet-workbench_code_delivery_2026-06-17\fg-meet-workbench_code_delivery_2026-06-17
\end{code}

在 MATLAB 命令行中执行：

\begin{code}
cd('C:\Users\You Huiwen\Desktop\fg-meet-workbench_code_delivery_2026-06-17\fg-meet-workbench_code_delivery_2026-06-17')
\end{code}

\section{步骤一：初始化路径}

进入交付包根目录后，先执行：

\begin{code}
setup_paths
\end{code}

该命令会加入以下路径：

\begin{table}[h]
\centering
\begin{tabular}{ll}
\toprule
路径 & 作用 \\
\midrule
\texttt{materials/} & 功能梯度分布、材料混合律和材料层参数函数 \\
\texttt{tools/} & 输入文件生成和辅助工具 \\
\texttt{matlab/} & 统一 MATLAB 求解入口 \\
\texttt{matlab/meet-fem-core/} & 实验室原有有限元核心程序 \\
\bottomrule
\end{tabular}
\end{table}

若未先执行 \texttt{setup\_paths}，MATLAB 可能无法找到求解器核心函数或材料辅助函数。

\section{步骤二：运行快速验证算例}

推荐先运行阶段 1 静力验证算例，用于确认路径、输入文件读取、矩阵装配和求解流程均正常：

\begin{code}
run('run_phase1_static_elastic.m')
\end{code}

该脚本默认读取：

\begin{code}
cases/Thermal_CFFF_U_Vf0.6-30x30-10layer.txt
\end{code}

并执行力载荷静力计算。正常运行后会生成：

\begin{code}
output/phase1_static_elastic_summary.csv
\end{code}

\section{步骤三：指定单个工况运行}

若需要手动指定某一个 \texttt{txt} 输入文件，可以执行：

\begin{code}
paths = setup_paths;

caseFile = fullfile(paths.cases, ...
    'Thermal_CFFF_U_Vf0.6-30x30-10layer.txt');

result = run_meet_static(caseFile, 'elastic', ...
    'OutTag', 'U_Vf06_check');
\end{code}

参数含义如下：

\begin{table}[h]
\centering
\begin{tabular}{ll}
\toprule
参数 & 含义 \\
\midrule
\texttt{caseFile} & 待读取的 MEET \texttt{txt} 输入文件 \\
\texttt{elastic} & 力载荷工况 \\
\texttt{OutTag} & 输出文件标签，用于区分不同运行结果 \\
\bottomrule
\end{tabular}
\end{table}

\section{步骤四：运行不同载荷类型}

同一个 \texttt{txt} 输入文件可以用于不同载荷类型。

\subsection{力载荷}

\begin{code}
result_elastic = run_meet_static(caseFile, 'elastic', ...
    'OutTag', 'U_Vf06_elastic');
\end{code}

\subsection{电载荷}

\begin{code}
result_electro = run_meet_static(caseFile, 'electro', ...
    'OutTag', 'U_Vf06_electro');
\end{code}

\subsection{磁载荷}

\begin{code}
result_magneto = run_meet_static(caseFile, 'magneto', ...
    'OutTag', 'U_Vf06_magneto');
\end{code}

输出结果通常包括中心挠度、层平均温差、温差跨度等指标。

\section{实际调用链路}

正确执行时，程序调用链路如下：

\begin{code}
run_meet_static(caseFile, loadCase)
    -> Main_FOSDLIN851T5MEET_V4(InputFile, ...)
        -> SF_GetInputDataMEEP(InputFile)
\end{code}

因此，实际矩阵装配和求解仍然使用实验室原有有限元核心程序；节点、单元、材料、边界条件等信息仍然来自 \texttt{txt} 输入文件。

\section{重新生成输入文件}

如需重新生成或刷新功能梯度输入文件，可在交付包根目录执行：

\begin{code}
python tools/generate_cases.py
\end{code}

该脚本会生成：

\begin{code}
cases/manifest_full.csv
cases/pilot_cases.csv
cases/Thermal_CFFF_*.txt
\end{code}

生成后的 \texttt{txt} 文件即为后续有限元求解的输入文件。

\section{材料辅助函数的用途}

\texttt{build\_layer\_materials.m} 是材料层参数生成函数，不是主求解程序。正确调用方式示例如下：

\begin{code}
setup_paths
layers = build_layer_materials(10, 6e-3, 0.6, 'U');
\end{code}

该函数只生成 10 层材料参数，不执行有限元装配和求解。因此不能在 MATLAB 命令行中零参数直接执行：

\begin{code}
build_layer_materials
\end{code}

否则 MATLAB 会报“输入参数数目不足”。该错误仅说明辅助函数调用方式错误，不能作为主程序无法运行的依据。

\section{推荐完整执行顺序}

推荐先按以下顺序进行最小验证：

\begin{code}
cd('交付包根目录')
setup_paths
run('run_phase1_static_elastic.m')
\end{code}

若阶段 1 算例正常完成，再运行指定工况：

\begin{code}
paths = setup_paths;
caseFile = fullfile(paths.cases, ...
    'Thermal_CFFF_U_Vf0.6-30x30-10layer.txt');

result_elastic = run_meet_static(caseFile, 'elastic');
result_electro = run_meet_static(caseFile, 'electro');
result_magneto = run_meet_static(caseFile, 'magneto');
\end{code}

\section{正常运行判断标准}

若程序正常运行，应至少出现以下结果之一：

\begin{itemize}
  \item MATLAB 命令行输出有限元装配和求解过程；
  \item 生成 \texttt{output/phase1\_static\_elastic\_summary.csv}；
  \item 生成对应的 \texttt{output/*.csv} 或 \texttt{output/*.mat}；
  \item MATLAB 工作区中出现 \texttt{result.wCenter\_mm} 等结果字段。
\end{itemize}

\section{执行注意事项}

\begin{enumerate}
  \item 不应直接进入 \texttt{materials} 文件夹运行单个辅助函数。
  \item 不应将 \texttt{build\_layer\_materials.m} 当作主程序。
  \item 正确入口应从交付包根目录执行 \texttt{setup\_paths}。
  \item 主求解应通过 \texttt{run\_phase1\_static\_elastic.m} 或 \texttt{run\_meet\_static} 执行。
  \item 输入数据来自 \texttt{cases/*.txt} 或实验室原始 \texttt{InputFile/*.txt} 文件。
\end{enumerate}

\end{document}
